\documentclass[11pt]{article}
\usepackage[margin=1in]{geometry}
\usepackage{amsmath, amssymb}
\usepackage{tikz}
\usetikzlibrary{arrows.meta, positioning}
\usepackage{booktabs}
\usepackage{microtype}
\usepackage{xcolor}

\tikzset{
  ver/.style ={draw, rounded corners, fill=black!6, font=\small, align=center, inner sep=5pt},
  verB/.style={draw, rounded corners, fill=blue!8, font=\small, align=center, inner sep=5pt},
  verR/.style={draw, rounded corners, fill=red!8, font=\small, align=center, inner sep=5pt},
  e/.style   ={-{Stealth[length=2.2mm]}, line width=0.7pt},
  lab/.style ={font=\scriptsize\itshape, text=black!60},
}

\title{\bfseries What must be remembered about a deleted character?\\[2pt]
\large v2: positions are testimony; only ties force retention\\
(working note --- throwaway)}
\author{Sal / FP Launchpad --- KC + Claude}
\date{2026-07-14}

\begin{document}
\maketitle

\paragraph{The question.} When a character is deleted from a replicated
list, what must the system remember about it, and for how long? The
requirement throughout (``the contract''): once any user has seen
character $u$ displayed before character $v$, no user may ever see $v$
before $u$; and replicas with the same events must agree.

\paragraph{Correction from v1.} Version 1 argued that retention of a
dead timestamp is forced whenever a deleted character has living
descendants. That is true for designs whose merge \emph{re-renders
positions from identity keys}, but KC's proposal --- keep every node's
absolute range untouched through deletes and merges --- refutes the
general claim: on v1's own example it produces the right answer with
\emph{nothing} retained. This version records the corrected boundary,
machine-checked: \textbf{distinct positions are permanent testimony;
retention is forced only at ties.}

\paragraph{Encoding.} Every node stores its full absolute range
$Q = (lo, hi) \subseteq (0,1)$. Insert takes the lower half of the
remaining gap inside the parent's absolute range (first child of the
root: $(0,\tfrac12)$). Display: depth-first, children by descending
$lo$, equal $lo$ broken by newer timestamp first. Characters are named
by their timestamps (larger = later). Under KC's proposal, delete
removes the record and re-parents the children; \emph{no number ever
changes}.

\section{The four-event example --- and why it forces nothing}

\begin{center}
\begin{tabular}{@{}cll@{}}
\toprule
step & action & screen \\
\midrule
1 & type $d$ (time 1) & $d$ \\
2 & type $a$ (time 4), same position & $a\ d$ \\
  & \emph{--- both replicas sync ---} & $a\ d$ \\
3 & \textbf{Alice}: type $x$ (time 6) \emph{after} $d$ & $a\ d\ x$
    \quad$\Rightarrow$ \textbf{Alice has seen $a$ before $x$} \\
4 & \textbf{Bob}: delete $d$ & $a$ \\
5 & \emph{merge} $M_1$ (LCA $=$ the synced version) & ? \\
\bottomrule
\end{tabular}
\end{center}

\noindent Ranges: $d = (0,\tfrac{16}{32})$, $a =
(\tfrac{16}{32},\tfrac{24}{32})$, $x = (0,\tfrac{8}{32})$ (lower half
\emph{of $d$}). Four designs at $M_1$:

\begin{center}\small
\begin{tabular}{@{}llll@{}}
\toprule
design & what survives $d$'s death & $M_1$ & verdict \\
\midrule
RGA & $d$ kept as a tombstone & $a\ x$ & \checkmark \\
always keep the dead time & stamp $(1)$ on $x$ & $a\ x$ & \checkmark \\
keep only if $d$ conflicted & nothing ($d$ never conflicted) & $x\ a$ &
  $\boldsymbol\times$ flips Alice's $(a,x)$ \\
\textbf{keep the ranges (KC)} & \textbf{nothing --- $x$ stays at
  $(0,\tfrac{8}{32})$} & $a\ x$ & \checkmark \\
\bottomrule
\end{tabular}
\end{center}

\noindent The keep-range row is the correction: $x$'s position
$(0,\tfrac{8}{32})$, below $a$'s $(\tfrac{16}{32},\tfrac{24}{32})$,
\emph{is} $d$'s testimony --- carried geometrically, permanently, for
free. The conditional design fails here not because information must be
retained but because it re-renders from keys and then forgets the key.
This example refutes conditional forgetting in re-render designs; it
does \emph{not} force retention on position-keeping designs.

\section{Keep-the-range survives everything the conditional designs
failed}

Machine-checked (regressions from \texttt{litmus/contest\_tree.py} run
against the keep-range design):

\begin{center}\small
\begin{tabular}{@{}lll@{}}
\toprule
countermodel & killed & keep-range \\
\midrule
4-event example (above) & conditional retention & survives: $a\ x$ \\
subordination escape (7 events) & contest-only stamps & survives:
  $(4,6)$ order preserved \\
retroactive subordination (6 events) & condition-on-current-level stamps
  & survives: $(3,2)$ order preserved \\
\bottomrule
\end{tabular}
\end{center}

\noindent The reason is uniform: in all three histories every deciding
comparison is between \emph{distinct} ranges, and kept ranges never
change --- so every verdict ever displayed is re-derivable from live
positions alone, at every replica, forever.

\section{The one failing shape: ties}

Concurrent characters typed at the same position from identical states
mint \emph{identical} ranges --- position carries zero information
between them. Machine trace (the L25 shape, three branches):

\begin{center}\small
\begin{tabular}{@{}ll@{}}
\toprule
step & state \\
\midrule
$R_0$: type $6$ at root & $6 = (0,\tfrac{16}{32})$ \\
$R_M$ (concurrent): type $10$ at root & $10 = (0,\tfrac{16}{32})$
  \quad$\leftarrow$ \textbf{identical mint: a tie} \\
$M_1 = \mathrm{merge}$: tie broken by time & display $[10,\ 6]$ \\
type $16$, $22$ under $6$ & $16 = (0,\tfrac{8}{32})$,\ $22 =
  (\tfrac{8}{32},\tfrac{12}{32})$, both $\subset (0,\tfrac{16}{32})$ \\
  & display $[10,\ 6,\ 22,\ 16]$ \quad($10$ before $22$ and $16$) \\
$R_D$ (from $R_0$): delete $6$ & \\
final merge, ranges kept & display $[\,\mathbf{22},\ \mathbf{16},\
  \mathbf{10}\,]$ \quad \textbf{both pairs flipped} \\
\bottomrule
\end{tabular}
\end{center}

\noindent The mechanism: $6$'s range \emph{equals} $10$'s, so ``inside
$6$'' is numerically also ``inside $10$'' --- $22$'s kept range has
$lo = \tfrac{8}{32} > 0 = 10$'s $lo$, and the position sort puts $22$
\emph{above} $10$. The verdict ``$10$ before $6$'' was a fact about the
\emph{pair of timestamps}, decided at $M_1$; the heirs' sub-ranges nest
inside \emph{both} rivals' equal ranges and cannot stay subordinate by
geometry once their side of the tie dies. At random-DAG scale keep-range
fails 52/120 executions --- this shape is the conjectured root of all of
them (minimization pending).

\section{The sibling-chain connection}

The Q-ranges and the sibling links are two encodings of one structure,
and the tie analysis says exactly where they coincide and where they
part.

\begin{center}\small
\begin{tabular}{@{}ll@{}}
\toprule
sibling-link view & Q-range view \\
\midrule
node's birth pointers: (parent, displaced sibling) & node's mint range \\
distinct chain positions & distinct ranges \\
\textbf{two nodes displacing the same head} & \textbf{identical mint
  ranges --- a tie} \\
splicing a dead node out of the chain & deleting a record, ranges kept \\
recording the spliced link through a dead node & the tie stamp / dig \\
\bottomrule
\end{tabular}
\end{center}

\noindent The carve is deterministic: same parent and same displaced
head $\Rightarrow$ the same slice. So \emph{tie-ness in ranges is
exactly equality of the birth pointers} --- a tie is two nodes claiming
the same chain position, and the range encoding is the sibling chain
\emph{arithmetized}. Both encodings order distinct chain positions for
free and both run out of positional information at the same place: a
shared chain position, where only timestamps can decide. And the
retention question is the same question in both wardrobes: keeping the
tie verdict for a dead node's heirs \emph{is} keeping the spliced-out
link through the dead node.

The two encodings genuinely differ in one behavior (machine-checked,
\texttt{litmus/sibling\_tree.py}): what happens to the
\emph{displacers} of a tied node. Branch $A$ stacks $2$ then $5$ over
$1$; branch $B$ concurrently stacks $4$ over $1$ --- so $2$ and $4$ tie
(both displaced $1$), and $5$ displaced $2$. The range read sorts the
level globally: $[5, 4, 2, 1]$ ($5$ is newest, floats to the top). The
sibling read keeps $5$ glued above the $2$ it displaced: $[4, 5, 2, 1]$.
Sequentially the encodings agree on 300/300 random histories; they part
exactly at ties. Run-gluing is the Fugue-direction behavior --- a
different datatype, not a different encoding of this one.

\section{The sharpened claim}

\begin{center}
\fbox{\parbox{0.92\textwidth}{\textbf{Retention is forced only at
ties.} Distinct concurrent mints are ordered by permanent geometry ---
equivalently, by their positions in the sibling chains. Identical
concurrent mints (shared chain positions) are ordered by a timestamp
verdict that some state must carry for as long as either side has
living descendants --- as a tombstone (RGA), a stamp (always-keep), a
retained spliced link (sibling chains), or any equivalent.}}
\end{center}

\medskip
\noindent This re-opens a door the conditional-retention refutations
appeared to close. Their principle was: \emph{the decision whether heirs
inherit a dead rank must be a function of immutable data} --- ``was it
contested'' and ``was it subordinated'' are mutable, concurrently
accruing knowledge, and both conditions fell to countermodels. But
\textbf{tie-ness is immutable}: mint ranges (equivalently, birth
pointers) are never rewritten, so ``$u$ and $v$ claimed the same
position'' is a fixed fact about the pair. A design that keeps ranges
verbatim and stamps heirs \emph{only when their dead ancestor was tied}
is therefore a live candidate, not a refuted one. Its retention would be
$O(\text{tie families})$: zero on a purely sequential million-character
history, growing only with genuine same-position races.

\paragraph{Pre-registered risks (the machine decides; not yet run).}
\emph{(i) Posthumous ties}: a later insert can coincidentally mint a
dead node's exact range (machine-observed) --- though such pairs were
never co-displayed, so their order may be legitimately free. \emph{(ii)
Tied orphans}: a tie whose dead side was born and deleted inside one
branch, invisible to the merge that meets the survivor. \emph{(iii)
Mixed-comparator transitivity}: position where distinct, stamps at ties
--- the combined relation must remain a total order. \emph{(iv) Nested
tie families}: heirs of both sides of a tie meeting each other.

\section{Where the LCA helps --- scoped honestly}

The MRDT merge sees three states; ``dead in a branch, live in the LCA''
is visible for free, timestamp included --- power that two-way-merge
CRDTs lack. v1 argued this covers only the first merge after a death,
with retention forced one merge later. That argument stands \emph{only
for re-render designs}: keep-the-range needs no consultation at all for
distinct-range deaths, and for tie deaths the merge that processes the
death typically still holds the tied record in an input (in the trace
above, the final merge sees $6$ live in one input). The honest residue
is exactly risk \emph{(ii)}: ties whose dead side no input remembers.

\paragraph{Status.} Machine-checked: remembering nothing at all breaks
--- at ties (8 refuted designs); conditional retention on mutable
conditions breaks (2 refuted rules); always-remembering works and equals
the tombstoned RGA ($O(\text{history})$ worst case); keep-the-range
works everywhere except ties (52/120 at scale). Open: tie-only retention
(the candidate above), and stability-gated discard of whatever is
retained.

\paragraph{Provenance.} \texttt{whiteboard/litmus/}:
\texttt{contest\_tree.py} (designs and countermodel regressions, and the
always-keep control's lockstep equivalence with the tombstoned RGA),
\texttt{sibling\_tree.py} (the encoding-divergence probe). The
keep-range runs reproduce from the session transcript of 2026-07-14; all
table values machine-printed.

\end{document}
